\documentclass[12pt, onecolumn]{IEEEtran}
\usepackage[brazilian]{babel}
\usepackage[utf8]{inputenc}
\usepackage{amsmath}
\usepackage{mathtools}
\usepackage{indentfirst}
\usepackage{graphicx}
\usepackage{float}
\usepackage{booktabs}
\usepackage{siunitx}
\usepackage{tabularx}
\usepackage{caption}
\usepackage{subcaption}


\title{Relatório Experimental 5 - Prática de Física dos Dispositivos Eletrônicos}
\author{Felipe Rodrigues Sobrinho (17/0141764)}
\date{14/06/2022}


\begin{document}

\maketitle

\section{Introdução}

O efeito hall é um fenômeno que acontece em um dispositivo semicondutor quando exposto a um campo magnético. Através da força de Lorentz, acontece um deslocamento no fluxo elétrico sobre a pastilha do semicondutor, gerando um novo potêncial elétrico. Tal deslocamento no fluxo elétrico é diretamente proporcional ao fluxo magnético sobre essa pastilha.

O sensor hall é um circuito integrado que é composto tanto pela pastilha semicondutora, quanto por um circuito de amplificação do sinal para que a corrente de saída seja suficiente para alimentar uma carga de menor impedância.

Esse relatório, então, descreverá o comportamento de um sensor hall analisado em laboraorio sob um crivo determinado pelo reteiro experimental da disciplina.

\section{Resultados}

\section{Conclusão}

\end{document}
